\section{Where the surplus constraints come from}
\label{sec:constraints}

The baseline showed that the constraints $p_u=p_v=0$ remove exactly the unstable band, and the curved model showed that they are first class. A first-class constraint
generates a gauge symmetry, so the question is which symmetry, and whether it is a natural one.

\begin{theorem}[A worldline gauge symmetry]
\label{thm:wlgauge}
For arbitrary $A_\mu(x)$ and flat $h$, with $C=A_\mu h^{\mu\nu}k_\nu$, the constraints satisfy $\{p_u,H\}=-Cp_v$, $\{p_v,H\}=Cp_u$, $\{p_u,p_v\}=0$ and $\{\ell,H\}=0$. The extended action
\[
\int d\tau\,\bigl(p\cdot\dot x-NH-\lambda_1p_u-\lambda_2p_v\bigr)
\]
is invariant under
\[
\delta u=\varepsilon_1(\tau),\qquad \delta v=\varepsilon_2(\tau),\qquad \delta\lambda_1=\dot\varepsilon_1+NC\varepsilon_2,\qquad \delta\lambda_2=\dot\varepsilon_2-NC\varepsilon_1,
\]
all other variables fixed: local translations along the normal plane, which are gauge symmetries although they are not isometries when $F\neq0$, the algebra closing with
structure functions \cite{HenneauxTeitelboim92}. In Dirac quantisation the physical states are the functions $\Psi(x)$ on spacetime, and on them $\hat H\Psi=-\tfrac12\Box_h\Psi$;
the potential $A$ drops out.
\end{theorem}

\begin{proposition}[Gauging the isometry is not enough]
\label{prop:isometryonly}
In polar normal coordinates $\ell=-p_\phi$ and $p_u^2+p_v^2=p_r^2+p_\phi^2/r^2$. Gauging only the rotation $\ell$, the isometry of the normal plane, leaves
$p_0^2=|\mathbf p|^2+m^2-p_r^2$: the radial surplus momentum keeps the unstable band. For fields, the rotation-invariant modes $J_0(qr)e^{-i\omega t+i\bk\cdot\mathbf x}$ obey
$\omega^2=|\bk|^2+m^2-q^2$.
\end{proposition}

The symmetry that removes the instability is therefore the full group of local translations along the normal plane, not its rotation subgroup.

\begin{corollary}[Weak emergence for scalar matter]
\label{cor:weak}
If the normal translations are gauged, the scalar sector reduces to the Klein--Gordon field on $(M^4,h)$: its energy is bounded below, its evolution is unitary, and
$[\phi(x),\phi(y)]=0$ at spacelike separation for the cone of $h$ \cite{WeinbergQFT}. In the free theory the zero modes $\psi$ and the surplus modes $\chi$ do not mix, since
translation invariance along the normal plane makes the mixed kernel $K_{\psi\chi}$ vanish, so the Schur complement is trivial and $K_{\rm eff}=K_{\psi\psi}$; the content is not in
the kernel but in whether $\chi$ is physical, and under the gauge principle it is not. The causal cone is that of $h$, the metric of the selected four-plane: it is inherited
from the reduction, not derived from the six-dimensional sum.
\end{corollary}

What the reduction leaves over is the geometry of the normal plane itself, and it is not healthy.

\begin{theorem}[Every normal-plane vector is a ghost]
\label{thm:ghostvec}
For $ds^2=h+\sigma\bigl[(du+B^u)^2+(dv+B^v)^2\bigr]$ with flat $h$, $B^u=b_1x^1dx^2$ and $B^v=b_2x^0dx^3$, one finds $R_6=-\tfrac\sigma4\sum F_{\mu\nu}F^{\mu\nu}$, and for the
rotation ansatz $B^u=vA$, $B^v=-uA$, $R_6=-\sigma\,r^2F^2/4$, which at $\sigma=+1$ is the ambient scalar curvature of \cref{thm:ein6}. The coefficient is odd in $\sigma$, so the
healthy sign is that of a normal plane of the sign of space, $\sigma=-1$, the ordinary Kaluza--Klein case. By \cref{thm:selector}(ii) the normal plane of a Lorentzian
four-plane in $\R^{3,3}$ always has the sign of time. Every Kaluza--Klein vector from the normal plane, from its translations or its rotation, therefore gravitates with the
wrong sign, and by \cref{thm:wlgauge} none of them couples to spinless matter.
\end{theorem}

The failure condition ``ghosts or negative spectral weights'' is therefore met by the normal-plane vectors, and the photon cannot come from them. If the Maxwell field
emerges, it must come from elsewhere: from the univectorial potential of the interior reading, or from the field-strength reading $\R^{3,3}=\Lambda^2\R^{1,3}$, where the
energy is the positive Hermitian form $|\bB|^2+|\bE|^2$ (\cref{sec:spin}).

\subsection{What the causal analysis establishes}

\Cref{tab:causal} collects the outcome against the failure conditions fixed at the start of \cref{sec:baseline}.

\begin{table}[t]
\centering
\small
\begin{tabular}{@{}>{\raggedright\arraybackslash}p{2.3cm}>{\raggedright\arraybackslash}p{6.6cm}>{\raggedright\arraybackslash}p{5.1cm}@{}}
\toprule
stage & result & verdict\\
\midrule
baseline & growing modes and no ground state; exactly the band the surplus constraints remove; compactification makes tachyons & obstruction confirmed, cure located\\
reduction & kinematics: selector $=J$, cone $=$ twisted squares at one constituent; dynamics: no $SO(3,3)$-invariant scalar mechanism; the positive route breaks the boosts &
kinematics exact; dynamical selection fails as posed\\
effective theory & scalar matter: no mixing, $K_{\rm eff}=K_{\psi\psi}$; normal-plane vectors are ghosts & failure condition met by the vectors\\
causality & the scalar sector is microcausal for the cone of $h$ & conditional on gauged translations; inherited\\
Maxwell & normal-plane vectors excluded as photons & blocked on this route\\
\bottomrule
\end{tabular}
\caption{The causal analysis against the failure conditions fixed in advance (\cref{sec:baseline,sec:selector,sec:typesel,sec:constraints}).}
\label{tab:causal}
\end{table}

Weak emergence holds for scalar matter, conditionally on a gauge principle, local translations along the normal plane, that is consistent but chosen rather than derived.
Strong emergence through a scalar order parameter is excluded. What survives is a sharper statement of the question. The selector is the complex structure $J$, and every
construction here that produced a healthy Lorentzian sector also used one extra datum: a positive inner product, which is the transpose of the realisation, equivalently
the temporal triple. Within these computations the totality of $SO(3,3)$-invariant dynamics fixes at most the type of the structure, never the instance, and it fixes even
the type only imperfectly. The datum that singles out time, wherever things worked, is the positivity of an inner product; in quantum mechanics it is the positivity of
the Hilbert-space norm. That is where a derivation of causal order from the split geometry would have to look.

Several routes remain open. A gauge mechanism for type selection, an $SO(3,3)$ analogue of the MacDowell--Mansouri and Stelle--West breaking
\cite{MacDowellMansouri77,StelleWest80}, $SO(3,3)\to SO(2)\times SO(1,3)$, would make the would-be ghosts gauge, solder the eight broken directions to the base and remove the six
flat ones. A principle that makes the normal translations gauge, rather than adding the multipliers by hand, would turn weak emergence from a conditional statement into a
derived one. At three constituents, a $\mathrm{Spin}(9,9)$-covariant construction of the cone of $E$, for which pure spinors are the natural candidates, or a proof that
none exists without the positive inner product, would settle whether the cone survives the tower.
